% 縦組での自己検査。既存の gckanbun-edge-test-tate.tex と同じ組み方
% （jlreq の tate クラスオプション）に合わせる。
%
% 当初の予測「箱ローカル座標は横組と同じ値になる」は外れた。実測では
% 縦組は LuaTeX-ja がグリフをもう1階層深い箱に包み、その包み箱が一律の
% オフセットを持つ（横組比で x に -0.5zw、y にも定数）。箱自体も
% h = d = 0.5zw の縦組メトリックになる。
%
% 一方でグリフ間の距離は横組と一致する。手計算で期待値を書けるのは
% 相対座標の方なので、縦組の検査は相対座標で行う。
\documentclass[luatex,fontsize=9pt,paper=a4,tate]{jlreq}
\usepackage{luatexja,lltjext}

\directlua{ NPT = dofile("gckanbun-nodepos-selftest.lua") }

\newsavebox{\Probe}

\newcommand{\Check}[3]{%
  \sbox{\Probe}{#3}%
  \directlua{ NPT.check(\number\Probe, "#1", "#2", true) }%
}

\begin{document}

% (1) 素直な連結。1zw の間隔が縦組でも保たれること
\Check{tate-plain}{U+5929/1->U+5730/1=589824,0}{\hbox{天地}}

% (2) 負のカーン。1zw - 0.5zw = 0.5zw
\Check{tate-negative-kern}{U+5929/1->U+5730/1=294912,0}{\hbox{天\kern-0.5\zw 地}}

% (3) \hbox to の伸長。自然幅 2zw を 3zw へ伸ばすので間隔は 2zw
\Check{tate-hbox-to-stretch}{U+5929/1->U+5730/1=1179648,0}{\hbox to 3\zw{天\hfil 地}}

% (4) 幅ゼロ右揃え箱の中の kern（特殊返り点が使う形）。
%     \makebox は 天 の直後、すなわち x=1zw に置かれる。中身の自然幅は
%     地(1zw) + kern(0.5zw) = 1.5zw で、右揃えなので 1zw - 1.5zw = -0.5zw から
%     始まる。よって 天 からの差は -0.5zw = -294912。
\Check{tate-makebox-kern}{U+5929/1->U+5730/1=-294912,0}%
  {\hbox{天\makebox[0pt][r]{地\kern0.5\zw}}}

\directlua{
  if NPT.report() > 0 then
    tex.error("nodepos tate selftest failed")
  end
}

縦組自己検査完了。

\end{document}
